%abstract here
Intraoral scanning is commonly performed manually by trained dental personnel, making the process costly, time-consuming, and dependent on operator experience. At the same time, automating this task is challenging, because a robot must move safely around complex dental geometry while capturing sufficient data without causing any harm to the patient. This project investigates whether existing robotics software, scanning, and computer vision technologies can be combined into an automated system for intraoral dental scanning and modeling.
\newline\newline
The project aims to create a pipeline that encompasses the process of scanning dental geometry and creating a 3D model of said dental geometry into a continuous integration, to evaluate if automation of this process is feasible. The concrete setup is roughly as follows: A simulation of an UR5e collaborative robot, equipped with a camera toolhead creates synthetic images of dental geometry using preexisting datasets. This is then operated through a multitude of preexisting components and libraries, that are chained together in order to recreate the initial jaw model using stereo reconstruction and registration.
\newline\newline
The scope of the project includes a pybullet simulation that renders multiple image pairs belonging to specific jawsides using blender, an analysis of different methods for achieving a 3D reconstruction from the rectified image pair and exporting it as a point cloud in .ply format, a xxx... and a prototype of a touchscreen user interface, built with Qt Creator that functions as a control panel, tracking a patient's head position and orientation through a camera and estimating the target coordinates a robotic arm would need to reach a selected tooth during a dental scanning procedure.
\newline\newline
Evaluation is currently done by using external datasets on their respective components, as they weren't yet developed far enough, to test the complete pipeline. However, as they have been designed to connect together, with more work, a finished pipeline could be evaluated by comparing the input 3D model to the reconstructed \& registrated 3D model.
\newline\newline
The resulting system may provide a basis for further investigating automated intraoral scanning under controlled conditions. The simulated image generation for example could be improved by using gazebo, with a much better ROS integration instead of pybullet, which in turn could allow using a real ur5e for image generation. Similarly, The head position and orientation estimation could be improved by replacing the current 2D landmark-based approximation with a real 3D calculation, using a calibrated camera matrix together with a method such as \texttt{solvePnP}. And in the future, the results of this project may be used to create a specialized medical robot, designed to be certified to work with real humans, thus making trained personnel available for other tasks.
